
\medskip
\noindent\SGARefTarget{sga3:iv:proposition:4:5:6}\textbf{Proposition 4.5.6.}
\textit{Let \(F\) and \(G\) be two sheaves over \(S\), and let
\(f:F\to G\) be an \(S\)-morphism. The following conditions on \(f\) are
equivalent:}
\begin{enumerate}
  \item[(i)]\SGARefTarget{sga3:iv:proposition:4:5:6:item:i:01} \textit{\(f\) is a monomorphism (respectively, an
    epimorphism, respectively, an isomorphism) in
    \(\widetilde{\mathscr C}\);}
  \item[(ii)]\SGARefTarget{sga3:iv:proposition:4:5:6:item:ii:01} \textit{\(f\) is a monomorphism (respectively, an
    epimorphism, respectively, an isomorphism) in
    \(\widetilde{\mathscr C}_{/S}=\widetilde{\mathscr C/S}\).}
\end{enumerate}

For monomorphisms and isomorphisms this is immediate, since it is a
question about presheaves. For epimorphisms it follows from the description
of epimorphisms as covering morphisms (\SGARefLink{sga3:iv:proposition:4:4:3}{4.4.3}) and from the fact that, by \SGARefLink{sga3:iv:subsection:4:5}{the
definition} of the induced topology, these are the same in
\(\widetilde{\mathscr C}\) and in \(\widetilde{\mathscr C/S}\).

\medskip
\noindent\SGARefTarget{sga3:iv:proposition:4:5:7}\textbf{Proposition 4.5.7.}
\textit{Let \(f:F\to G\) be a morphism of sheaves. The following conditions
are equivalent:}
\begin{enumerate}
  \item[(i)]\SGARefTarget{sga3:iv:proposition:4:5:7:item:i:01} \textit{\(f\) is a monomorphism (respectively, an
    epimorphism, respectively, an isomorphism);}
  \item[(ii)]\SGARefTarget{sga3:iv:proposition:4:5:7:item:ii:01} \textit{for every \(S\in\operatorname{Ob}(\mathscr C)\),
    the morphism \(f_S:F_S\to G_S\) is locally a monomorphism
    (respectively, an epimorphism, respectively, an isomorphism), that is:}
  \item[(iii)]\SGARefTarget{sga3:iv:proposition:4:5:7:item:iii:01} \textit{for every \(S\in\operatorname{Ob}(\mathscr C)\),
    the set of morphisms \(T\to S\) for which
    \(F_T\to G_T\) is a monomorphism (respectively, an epimorphism,
    respectively, an isomorphism) is a refinement of \(S\).}
\end{enumerate}
\textit{If the given topology is defined by a pretopology \(\mathscr R\),
these conditions are also equivalent to:}
\begin{enumerate}
  \item[(iv)]\SGARefTarget{sga3:iv:proposition:4:5:7:item:iv:01} \textit{for every \(S\in\operatorname{Ob}(\mathscr C)\),
    there is a covering family
    \(\{S_i\to S\}\in\mathscr R(S)\) such that, for every \(i\),
    \(F_{S_i}\to G_{S_i}\) is a monomorphism (respectively, an
    epimorphism, respectively, an isomorphism).}
\end{enumerate}
\textit{If \(\mathscr C\) has a final object \(e\), it is enough in
conditions \SGARefLink{sga3:iv:proposition:4:5:7:item:ii:01}{(ii)}, \SGARefLink{sga3:iv:proposition:4:5:7:item:iii:01}{(iii)}, and \SGARefLink{sga3:iv:proposition:4:5:7:item:iv:01}{(iv)} to take \(S=e\).}

Clearly \SGARefLink{sga3:iv:proposition:4:5:7:item:ii:01}{(ii)}, \SGARefLink{sga3:iv:proposition:4:5:7:item:iii:01}{(iii)}, and \SGARefLink{sga3:iv:proposition:4:5:7:item:iv:01}{(iv)} are equivalent. To prove the equivalence of
\SGARefLink{sga3:iv:proposition:4:5:7:item:i:01}{(i)} and \SGARefLink{sga3:iv:proposition:4:5:7:item:ii:01}{(ii)}, as well as \SGARefLink{sga3:iv:proposition:4:5:7}{the assertion} concerning a final object, it is
enough to show that \SGARefLink{sga3:iv:proposition:4:5:7:item:ii:01}{(ii)} implies \SGARefLink{sga3:iv:proposition:4:5:7:item:i:01}{(i)} and that the properties in question
are stable under base change. We prove the latter point first. For
monomorphisms and isomorphisms it is immediate, since it is a question
about presheaves. For epimorphisms it follows from the fact that every
epimorphism of sheaves is universal (\SGARefLink{sga3:iv:proposition:4:4:3}{4.4.3}).

It remains to show that \SGARefLink{sga3:iv:proposition:4:5:7:item:ii:01}{(ii)} implies \SGARefLink{sga3:iv:proposition:4:5:7:item:i:01}{(i)}. Suppose first that
\(f_S:F_S\to G_S\) is locally a monomorphism (respectively, an
isomorphism). There is then a covering sieve \(C\) of \(S\) such that
\(f_T\) is a monomorphism (respectively, an isomorphism) for every
\(T\to S\) belonging to \(C\). Since an inverse limit of monomorphisms
(respectively, isomorphisms) is again one,
\(\operatorname{Hom}(C,f)\) is a monomorphism (respectively, an
isomorphism); see \SGARefLink{sga3:iv:proposition:4:1:4}{4.1.4}. The commutative diagram
\[
\begin{array}{ccc}
 \operatorname{Hom}(C,F)
 &\xrightarrow{\operatorname{Hom}(C,f)}&
 \operatorname{Hom}(C,G)
 \\
 \big\uparrow\wr && \big\uparrow\wr
 \\
 F(S)&\xrightarrow{f(S)}&G(S)
\end{array}
\]
therefore shows that \(f(S)\) is injective (respectively, bijective).

Finally, suppose that \(f:F\to G\) is locally an epimorphism, and let
\(H\subset G\) be its image. For every \(S\to G\), the sheaf
\(H\times_GS\) is the image of
\[
  f\times_GS:F\times_GS\longrightarrow S.
\]
To prove that \(f\) is an epimorphism, it is enough to prove that \(H\) is
a refinement of \(G\), or equivalently that \(H\times_GS\) is a
refinement of \(S\) for every \(S\). But this holds after each base change
\(T\to S\) belonging to a refinement of \(S\), since
\(f\times_GS\) is locally covering. Hence \(H\times_GS\) is itself a
refinement of \(S\), by Axiom \SGARefLink{sga3:iv:definition:4:2:1:item:t:2:01}{(T~2)}.

\medskip
\noindent\SGARefTarget{sga3:iv:corollary:4:5:8}\textbf{Corollary 4.5.8.}
\textit{Let \(F\) and \(G\) be two sheaves over \(S\), and let
\(f:F\to G\) be an \(S\)-morphism. Then \(f\) is a monomorphism,
respectively an epimorphism, respectively an isomorphism, if and only if it
is so locally on \(S\).}

\medskip
\noindent\SGARefTarget{sga3:iv:remark:4:5:9}\textbf{Remark 4.5.9.}
The proof of \SGARefLink{sga3:iv:proposition:4:5:7}{the proposition} shows that its assertion about monomorphisms
(respectively, isomorphisms) remains valid if one assumes only that \(F\)
is a separated presheaf (respectively, a sheaf) and \(G\) is an arbitrary
presheaf (respectively, a separated presheaf).

